% !TEX TS-program = xelatex
% !TEX encoding = UTF-8 Unicode
% !Mode:: "TeX:UTF-8"

\documentclass{resume}
\usepackage{zh_CN-Adobefonts_external} % Simplified Chinese Support using external fonts (./fonts/zh_CN-Adobe/)
%\usepackage{zh_CN-Adobefonts_internal} % Simplified Chinese Support using system fonts
\usepackage{linespacing_fix} % disable extra space before next section
\usepackage{cite}
\usepackage[svgnames]{xcolor}
\usepackage{hyperref}
\usepackage{fontawesome}
\usepackage{graphicx} 
\hypersetup{colorlinks,breaklinks,
            urlcolor=[rgb]{0.2,0.2,0.6},
            linkcolor=[rgb]{0.2,0.2,0.6}
}

\begin{document}
\pagenumbering{gobble} % suppress displaying page number


% \name{Ziyang Feng}
% \vspace{-0.5ex}
% \centerline{Software Engineer: Backend Dev. \& Big Data Dev.}
% \vspace{1ex}
% % {E-mail}{mobilephone}{homepage}
% % be careful of _ in emaill address
% % {E-mail}{mobilephone}
% % keep the last empty braces!
% \contactInfo{fuujiro@qq.com}{(+86) 155-2489-2259}{ \faGithub\ \url{https://github.com/fuujiro}}
% \vspace{-1.5ex}
 
% \section{\faGraduationCap\ Education}
% \vspace{-0.5ex}
% \datedsubsection{\textbf{Waseda University}}{Sept. 2020 -- Expected Jul. 2022}
% \vspace{-0.5ex}
% \textit{Master of Engineering} in Information, Production and Systems
% \vspace{-1ex}
% \datedsubsection{\textbf{Dalian University of Technology}}{Sept. 2016 -- Jul. 2020}
% \vspace{-0.5ex}
% \textit{Bachelor of Engineering} in Computer Science \& Technology
% \vspace{-0.6ex}

% \section{\faSitemap\ Intern Experience}
% \vspace{-0.5ex}
% \datedsubsection{\textbf{\includegraphics[scale=0.5]{img/alibabacloud.eps} Alibaba Cloud - Computing Platform - DataWorks}}{Apr. 2021 -- Present}
% %\role{\LaTeX, Maintainer}{Individual Projects} \includegraphics[]{img/alibabacloud.pdf}
% %The goal of this task was to detect and correct spelling errors on Chinese essays.
% \begin{itemize}
%   \item Using SofaArk framework and reflection mechanism to implement plug-in registration and reference; used Quartz task scheduling framework to implement heartbeat failover mechanism to ensure high availability.
%   \item Developed a series of checkers, such as violence scan, CU consumption detection, etc. to guarantee data quality; developed a series of data governance APIs such as health score, whitelist mechanism, etc.
% \end{itemize}
% \vspace{-2.5ex}

% \datedsubsection{\textbf{\includegraphics[scale=0.5]{img/tencentqq.pdf} Tencent Games - LightSpeed \& Quantum Studios}}{Dec. 2020 -- Apr. 2021}
% %\role{\LaTeX, Maintainer}{Individual Projects} \includegraphics[]{img/alibabacloud.pdf}
% %The goal of this task was to detect and correct spelling errors on Chinese essays.
% \begin{itemize}
%   \item With ZeroMQ, implemented an asynchronous and high-performance C/C++ service framework based on the Router-Dealer prototype, which provides asynchronous and concurrent support for game AI training.
%   \item Use Python to develop integrated tools, such as automatically pulling and matching the md5 value; Use deep learning inference frameworks NCNN and TNN for end-to-side acceleration and improve AI functions.
% \end{itemize}
% \vspace{-2.5ex}

% \datedsubsection{\textbf{\includegraphics[scale=0.5]{img/huawei.pdf} Huawei Technologies Co., Ltd. - CloudBU}}{Jun. 2019 -- Aug. 2019}
% %\role{\LaTeX, Maintainer}{Individual Projects}
% %The goal of this task was to detect and correct spelling errors on Chinese essays.
% \begin{itemize}
%   \item Use C to complete load balancing of the gateway, and the limiting algorithm (leaky bucket, token bucket).
%   \item Based on Python and Go to complete the IP script test tool, optimize the slow query of the client's IP.
% \end{itemize}
% \vspace{-1.5ex}

% \section{\faUsers\ Project Experience}
% \vspace{-0.5ex}
% \datedsubsection{\textbf{RPC Framework(\href{https://github.com/fuujiro/jiro-rpc-framework}{jiro-rpc-framework})}}{Sept. 2020 -- Oct. 2020}
% %\role{Summer Intern}{Manager: xxx}
% \begin{itemize}
%   \item Through dynamic proxy, with the Netty transmission architecture, implemented automatic registration with asynchronous non-blocking and heartbeat notification, implemented a variety of serialization methods.
%   \item Service registration and discovery based on Nacos realizes automatic service cancellation and load balancing strategies (random distribution, balanced weighted round-robin, consistent Hash).
% %  \item Optimized xxx 5\%
% \end{itemize}
% \vspace{-2.5ex}

% \datedsubsection{\textbf{Distributed Systems (\href{https://github.com/fuujiro/MIT6.824-Labs}{MIT6.824-Labs})}}{May. 2020 -- Oct. 2020}
% %\role{C, Python, Django, Linux}{Individual Projects, collaborated with xxx}
% %Brief introduction: xxx
% \begin{itemize}
%   \item Based on the MapReduce, implemented the master's distribution task and the worker's calculation task.
%   \item Understand the Raft protocol and the heartbeat mechanism, such as the election and tenure, log synchronization and compression, member changes, and strong consistency.
% %  \item xxx
% \end{itemize}
% \vspace{-1.5ex}

% % Reference Test
% %\datedsubsection{\textbf{Paper Title\cite{zaharia2012resilient}}}{May. 2015}
% %An xxx optimized for xxx\cite{verma2015large}
% %\begin{itemize}
% %  \item main contribution
% %\end{itemize}


% \section{\faCogs\ Skills}
% \begin{itemize}[parsep=0.5ex]
%   \item Understand the operating system's memory management, TCP/IP protocol stack, and IO multiplexing.
%   \item Skilled in Java, C++, Git and \LaTeX, familiar with network programming, had good programming style.
%   \item Experienced in Spring, SpringBoot, MyBatis and other back-end frameworks, and basic design patterns.
%   \item Experienced in Java concurrent programming, thread pool mechanism and the understanding of JUC library.
%   \item Understand JVM memory distribution, class loader mechanism, garbage collection algorithm.
%   \item Familiar with middleware such as Kafka, ElasticSearch and Thymeleaf template engine.
%   \item Familiar with Linux command shell and Docker \& virtual environments, and CI/CD Tools.
% \end{itemize}
% \vspace{-1.5ex}

% \section{\faTrophy\ Academic Competitions}
% \datedline{\textit{National University Student Innovation Project}~~Awarded in Robot-arm vision calibration}{Mar. 2019}
% \datedline{\textit{\nth{1} Prize}~~Awarded in Liaoning Province University Student Computer Application Competition}{Dec. 2018}
% \datedline{\textit{\nth{1} Prize}~~Awarded in National College Student Mathematics Competition Dalian Division}{Jul. 2017}

% \clearpage

\name{冯子扬 }
\centerline{求职意向: 后端开发 | 大数据研发}
% {E-mail}{mobilephone}{homepage}
% be careful of _ in emaill address
\vspace{1ex}
\contactInfo{fuujiro@qq.com}{(+86) 155-2489-2259}{\faGithub\ \url{https://github.com/fuujiro/} }

% \vspace{-1ex}
% \vspace{-1ex}

\vspace{-1ex}
 
\section{\faGraduationCap\  教育背景}
\datedsubsection{\textbf{早稻田大学}~~ \ 硕士, 信息系统 }{2020.09 -- 2022.09}

\datedsubsection{\textbf{大连理工大学}~~ \ 学士, 计算机科学与技术 }{2016.09 -- 2020.06}
\vspace{1ex}

%\vspace{-1ex}

\section{\faSitemap\ 工作经历}
\datedsubsection{\textbf{ 阿里云 - 计算平台DataWorks ~/~ 后端开发}}{2022.09 -- 至今为止}
\vspace{-0.5ex}
%\role{\LaTeX, Python}{个人项目}
\begin{onehalfspacing}
%优雅的 \LaTeX\ 简历模板, https://github.com/fuujiro/myResume
\begin{itemize}
  \item 开发了检查器以及其注册配置框架，每日触发检查数万次，建设高可用的failover心跳机制。
  \item 基于MaxCompute离线大数据引擎开发，搭建了完整的离线数仓分层加工框架，监测表、任务数据的特征与波动，提供针对存储、计算的治理方案和工具。
  \item 实现OLAP（Hologres）与OLTP（MySQL）的多数据源访问机制，针对性响应页面接口优化，使用Flink 链路来保证数据库一致性，以及提供监控方案和降级策略。
\end{itemize}
\end{onehalfspacing}

\section{\faUsers\ 项目经历}
\datedsubsection{\textbf{检查器} }{}
%\role{实习}{经理: 高富帅}
%xxx后端开发
\vspace{-0.5ex}
\begin{onehalfspacing}
\begin{itemize}
  \item 检查器是作用于数据治理的事前阶段，通过SQL解析，开发了20余个检查器（如高费用SQL、表及字段存在性\&一致性检查等）,对用户运行、提交、发布的代码判断拦截，日检查达万次量级
  \item 通过反射机制来实现检查器的注册和引用，框架也支持二方包机制的集团内的三方接入
  \item 建设了高可用的failover机制，实现各机器定时上报心跳，断线抢主和对失败任务二次分发
\end{itemize}
\end{onehalfspacing}

\vspace{-1.5ex}

\datedsubsection{\textbf{治理项 \& 离线数仓 \& OpenAPI}}{}
\vspace{-0.5ex}
\begin{onehalfspacing}
\begin{itemize}
  \item 治理项是作用于事中阶段，设计五大治理领域（研发、存储、计算、质量、安全），开发了60余个治理项（如上游依赖缺失、任务运行时间超长，持续导入一致等），日产出百万量级的治理项问题
  \item 建设了一套完整的大数据离线框架，从Mysql同步至MaxCompute离线大数据引擎，进行T+1的任务、表血缘链路追踪、定位、加工、收益计算、健康打分、回流的全链路加工流程
  \item 设计并开发一系列高可用、高质量、高响应的数据治理API接口（如列出治理项问题、检查器配置、任务优雅下线、加白名单等）
\end{itemize}
\end{onehalfspacing}

\datedsubsection{\textbf{数据质量 \& 治理工具}}{}
\vspace{-0.5ex}
\begin{onehalfspacing}
\begin{itemize}
  \item 为各治理项开发自动化工具，如一键补全上游依赖，分析表和任务血缘的差异实现自动化补全
  \item 优化页面响应，深度调优SQL性能优化，实现了MySQL和Hologres的多数据源访问方案，使用flink cdc来保证数据源的一致性，以及对应的监控和降级策略
  \item 提供专项自动化治理方案，如任务优雅下线（提供完善的下线链路，如延迟调度、暂定虚拟根节点、重命名产出表、一键回滚、下线API等），节点依赖切换等
  \item 使用观测者模式和策略模式完善治理通知，提供治理项在各维度（租户、空间、个人）下各产品（钉钉、微信、飞书、webhook）的报表和提醒通知
\end{itemize}
\end{onehalfspacing}

% Reference Test
%\datedsubsection{\textbf{Paper Title\cite{zaharia2012resilient}}}{May. 2015}
%An xxx optimized for xxx\cite{verma2015large}
%\begin{itemize}F
%  \item main contribution
%\end{itemize}

%\vspace{-1.5ex}
%\section{\faTrophy\ 技术栈\&专利}
%  \item Java（Spring Boot） | SQL | Hive | MaxCompute
%  \item MySQL | PostgresSQL | Redis | Kafka | 
%  \item AWS | Alibaba Cloud | Git | Maven | Junit | pRO

\section{\faCogs\ 技术栈}
% increase linespacing [parsep=0.5ex]
\begin{itemize}[parsep=0.5ex]
  \item 熟悉 Java并发编程，SpringBoot框架，多线程和线程池机制以及常用JUC类库的底层原理。
  \item 理解 JVM内存分布，类加载、类隔离机制，垃圾回收算法和常用GC调优策略。
  \item 理解 OLAP与OLTP数据库原理，熟悉 MySQL的存储引擎、事务及隔离级别、锁、索引
  \item AWS | Alibaba Cloud | Docker | Git | Maven | Junit | Prometheus | Grafana
\end{itemize}

%% Reference
%\newpage
%\bibliographystyle{IEEETran}
%\bibliography{mycite}
\end{document}
